% SPDX-License-Identifier: CC-BY-4.0
% Copyright (c) 2026 Martin Schröder <info@swedishembedded.com>

\section{The earlier series in detail}\label{app:series}

\subsection{S1: self-taught generations, judged first pooled and then per lineage}\label{sec:rsi1}

S1 is one registration with dated arms, each added before it ran; its exams were selected and
graded by the arm's own registered analysis script, without a separate examiner.\cl{RS01} Arms A and B
ran the closed-book loop of \cref{sec:documentation} for three generations and were judged on the
loop's own probes; re-graded under a neutral request, most of the gain they measured was adaptation
to the probing protocol and did not transfer.\cl{RS11} Arm H therefore built the exam described
above and measured every generation on it greedily under the neutral request. Its three replicates,
trained only on tasks their own models wrote, answered 38, 31 and 34 of the 889 items at
generation~1, 60, 63 and 49 at generation~2 and 67, 83 and 59 at generation~3 (base 0); pooled over the
replicates the steps were 103/0, 95/26 and 103/66 wins/losses ($p=0.00273$ at the third) with no
replicate falling, which met the registered pooled criterion.\cl{RS12} The same design in Python
(arm HP, 820 items, base 2) met its first two steps (49/7 and 83/8) and not the third (59/50,
$p=0.222$), where one replicate fell from 69 to 60.\cl{RS13} Arm HR added replay across generations
and a pooled retention bound and met the pooled criterion (66/4, 16/1, 25/3), but each step's gain came
from a single replicate, so it is not evidence that one lineage improved twice.\cl{RS14} Three arms
changed the update instead: self-imitation of the model's own verified reasoning (R) and policy
gradient with negative rewards at two update scales (P, P2). After three iterations they answered 21,
4 and 12 of the 889 items, and on six families none of them trained on their pass@1 went from 0.41
to 0.4479, 0.4408 to 0.4242 and 0.4419 to 0.4408: none met its criterion.\cl{RS15}

Because a pooled test cannot tell staggered single gains from one model improving repeatedly, arm HL
repeated arm H's recipe in four fresh lineages under a per-lineage criterion: a lineage succeeds when
all three of its steps hold, and the arm when two of four do. Its first run ended on the project disk's
inode quota before any lineage finished and was not evaluated.\cl{RS16} The rerun met the criterion in
1 of 4 lineages: r3 rose at every step (16/0; 14/5, $p=0.0318$; 33/12, $p=0.00123$), r2 rose twice and
fell at generation~3 (27/36), r0's later steps were not significant (35/25, $p=0.123$; 37/28,
$p=0.161$), and r1 accepted no round after generation~1. Pooled, the third step was 97/76
($p=0.0641$).\cl{RS17} S1 thus showed large first steps on items the base never answered and
third steps that did not hold within a lineage.

\subsection{S2: on-policy policy gradient}\label{sec:rsi2}

S2 asked whether replacing the update with on-policy policy gradient under powered gates makes the
later steps hold. Each generation trained on the current policy's own graded attempts at programs it
solves sometimes ($k=8$ attempts, group-standardised advantages, clip edges 0.2 and 0.28, learning
rate $6\times10^{-6}$ for 24 steps, LoRA rank 16), and the exam moved to the independent
examiner.\cl{RS10,RS18} At generation~1 the four lineages won 4, 2, 3 and 3 items against the base and
lost none ($p$ from 0.0625 to 0.25): no step held, so no lineage could succeed, and the examiner's
plan verdict reads not met.\cl{RS19} The examiner's overlap audit found that the policy's own
proposals had reproduced 1, 1 and 2 exam programs exactly in three lineages; S2 had no exclusion
and those items count as scored.\cl{RS19} Generation~2 could not be served, because every trained
candidate carried the same adapter card id; the criterion was already decided, so it was not rerun
after the fix.\cl{RS08} The recipe raised the learner's own dev set in its pilot but moved the exam by
two to four items.

\subsection{S3: training on the interpreter's answers to the policy's own programs}\label{sec:rsi3}

S3 returned to arm H's source of information, training on the interpreter's output for programs the
policy itself proposed, and replaced what had failed in S1: acceptance on a dev set powered for an
effect of 0.08, replay of earlier generations, proposals diversified per manual section, and the whole
learner behind Brain's API. A generation proposes up to 600 programs, predicts each greedily and
trains a LoRA continued from the champion (rank 16, learning rate $3\times10^{-5}$, 40 steps) on the
mispredicted programs, with the correctly predicted ones and earlier generations' records each drawn
at a share of 0.25; the examiner now excludes overlapping items.\cl{RS09,RS05}

Every lineage improved at generations~1 and~2; at generation~3 r2 and r3 improved again, r0 held
flat (29/28, $p=0.5$) and r1 fell (18/28) (\cref{tab:rsi}). Two of four lineages succeeded and the
plan verdict reads not met; pooled, the steps were 288/0, 199/60 and 142/92, and no exam item matched a
training program.\cl{RS21,RS22} The learner's records show the signal running out: over the four
lineages the champions mispredicted 2127 of 2400 proposals at generation~1, 1125 of 2228 at
generation~2 and 860 of 1942 at generation~3, so later generations trained mostly on what their
predecessor already knew, while each also continued the previous adapter.\cl{RS23} The learner's dev
gate accepted all 12 steps, and the examiner found 2 of them not to hold: r0's and r1's third steps,
the two smallest dev margins (13/8 and 7/5).\cl{RS24}


\begin{figure}[tbp]
\centering
\begin{tikzpicture}
\begin{groupplot}[paper, group style={group size=2 by 1, horizontal sep=12mm},
  width=0.48\linewidth, height=5.4cm, xmin=-0.25, xmax=3.25, ymin=0, ymax=260,
  xtick={0,1,2,3}, xticklabels={base,1,2,3}, xlabel={generation},
  ylabel={items right of 889 (exam subset)}]
\nextgroupplot[title={S3: continued adapter (not met)}]
\addplot[lineage r0] table[col sep=comma,x=gen,y=r0] {generated/rsi3_items.csv};
\addplot[lineage r1] table[col sep=comma,x=gen,y=r1] {generated/rsi3_items.csv};
\addplot[lineage r2] table[col sep=comma,x=gen,y=r2] {generated/rsi3_items.csv};
\addplot[lineage r3] table[col sep=comma,x=gen,y=r3] {generated/rsi3_items.csv};
\nextgroupplot[title={S4: retrained from the base (met)}, ylabel={},
  legend style={at={(-0.11,-0.3)}, anchor=north, legend columns=4, /tikz/every even column/.append style={column sep=4mm}}]
\addplot[lineage r0] table[col sep=comma,x=gen,y=r0] {generated/rsi4_items.csv};
\addplot[lineage r1] table[col sep=comma,x=gen,y=r1] {generated/rsi4_items.csv};
\addplot[lineage r2] table[col sep=comma,x=gen,y=r2] {generated/rsi4_items.csv};
\addplot[lineage r3] table[col sep=comma,x=gen,y=r3] {generated/rsi4_items.csv};
\legend{lineage r0, lineage r1, lineage r2, lineage r3}
\end{groupplot}
\end{tikzpicture}
\caption{Items right on the held-out exam per generation and lineage, as the examiner measured them; the
exam holds items the base did not answer when it was selected, so these are counts on that hard
subset, not an accuracy on Lua programs. In S3 two lineages stop rising at generation~3 (r0 flat, r1 down); in S4
every lineage rises at every generation. The base level differs between the panels because the base
was served at a different KV tier (\cref{tab:rsi}).}\label{fig:rsi-items}
\end{figure}

\paragraph{Serving configuration.} The examiner served S4's models with an f32 KV cache, at which the base
answers 12 exam items; S3's stored baseline was served at the build's default tier, at which it answered
none, so S3's and S4's levels are not comparable, only their steps are.\cl{RS27}
